%%%%%%%%%%%%%%%%%%%%%%%%%%%%%%%%%%%%%%%%%%%%%%%%%%%%%%%%%%%%%%%%%%%%%%%%%%%%%%%%
% For the use of University of Amsterdam
% Master of AI thesis
%%%%%%%%%%%%%%%%%%%%%%%%%%%%%%%%%%%%%%%%%%%%%%%%%%%%%%%%%%%%%%%%%%%%%%%%%%%%%%%%

% NOTE: the [license] option is intentionally NOT set here. Adding a license
% (CC-BY-NC by default) depends on media/titlepage/by-nc.png and should only be
% enabled after confirming the IP of the work with your supervisors. To enable
% it, write \documentclass[license]{mscaithesis}.
\documentclass{mscaithesis}

\usepackage{mwe}  % provides example-image for the figure placeholders

\usepackage{amssymb} % Math packages are (on purpose) not included in the class document
\usepackage{amsfonts}
\usepackage{amsmath}
\usepackage{amsthm}

\usepackage{natbib}

%%%%%%%%%%%%%%%%%%%%%%%%%%%%%%%%%%%%%%%%%%%%%%%%%%%%%%%%%%%%%%%%%%%%%%%%%%%%%%%%
% DOCUMENT METADATA
%%%%%%%%%%%%%%%%%%%%%%%%%%%%%%%%%%%%%%%%%%%%%%%%%%%%%%%%%%%%%%%%%%%%%%%%%%%%%%%%

\title{Semantic Scene Completion with Latent\\ Modeling of Scene-Level 3D Gaussians}

% Date on which your thesis is submitted (TODO: set the real submission date)
\date{2026/06/30}

\author{Mohammad Hafeez \textsc{Khan}}
\studentnumber{TODO: 15549089 }      % TODO: fill in
\period{2025/11 -- 2026/06}               % TODO: confirm the official project period
\credits{36 ECTS}
\examiner{TODO: assigned examiner}        % TODO: fill from Datanose
\dailysupervisor{Yue Li}
\uvasupervisor{Dr.\ Martin R.\ Oswald}    % UvA supervisor (the ``other reader'')
\addsupervisor{}                          % optional; leave empty if not applicable

\begin{document}
\maketitle

\frontmatter
        \chapter*{Abstract}
                \begin{adjustwidth}{4cm}{4cm}
                        3D Gaussian Splatting has become a standard representation for photorealistic scene reconstruction. Generating and completing whole scenes in this representation first requires compressing a scene of many unordered Gaussian primitives into a structured latent space. This thesis studies a semantically aware, scalable, scene-level Gaussian-splatting autoencoder as the first stage of a two-stage generative pipeline. We identify the central obstacle to scaling such an autoencoder as the non-identifiability of the per-primitive parameters: a fitted scene contains large amounts of gauge freedom, in the ordering of the set, in the quaternion encoding of orientation, and in the placement of primitives within a region, which a naive reconstruction objective cannot learn. We address this by making the reconstruction target identifiable and the latent compositional, through a fixed-frame Hilbert serialization, permutation-invariant and gauge-invariant reconstruction objectives, and a spatially local per-token encoder, with contrastive and disentanglement objectives that give the latent semantic and layout structure. We show that these changes are what let the autoencoder scale from a few hundred scenes toward several thousand, analyze which choices matter most, and specify the second-stage generator for scene generation and completion that the representation is designed to support.
                \end{adjustwidth}

\tableofcontents

\listoffigures

\listoftables

\setlength{\parskip}{5pt}

\chapter{List of Symbols}

The following notation is used throughout the thesis.

\begin{table}[h]
  \centering
  \begin{tabular}{ll}
    $\mathcal{G} = \{g_i\}_{i=1}^{N}$ & a scene as a set of $N$ Gaussian primitives \\
    $g_i = (\mu_i, c_i, \alpha_i, s_i, q_i)$ & a single Gaussian primitive \\
    $\mu_i \in \mathbb{R}^3$ & primitive position \\
    $c_i \in \mathbb{R}^3$ & primitive (diffuse) colour \\
    $\alpha_i \in (0,1)$ & primitive opacity \\
    $s_i \in \mathbb{R}^3_{+}$ & primitive per-axis scale \\
    $q_i \in \mathbb{S}^3$ & primitive orientation (unit quaternion) \\
    $\Sigma_i = R(q_i)\,\mathrm{diag}(s_i^2)\,R(q_i)^\top$ & primitive covariance \\
    $R(q)$ & rotation matrix of a quaternion $q$ \\
    $N$ & number of primitives sampled per scene \\
    $E_\phi,\ D_\theta$ & encoder and decoder \\
    $Z \in \mathbb{R}^{T \times d}$ & latent of $T$ tokens of width $d$ \\
    $T = 512$ & number of latent tokens \\
    $d = 32$ & per-token latent dimension \\
    $z_s,\ z_g$ & layout (semantic) and geometry latent components \\
    $\beta$ & weight of the KL regulariser \\
    $K$ & number of Hilbert blocks (equal to $T$ in this work) \\
    $g = \lceil N/K \rceil$ & Gaussians per block \\
    $\omega$ & half-width of the local encoder window (in blocks) \\
    $a_{\tau(i)}$ & anchor of the token $\tau(i)$ owning primitive $i$ \\
    $o_i,\ \rho$ & raw position offset and learnable offset scale \\
    $h(i)$ & Hilbert-curve key of primitive $i$ \\
    $v_\theta(Z_t, t)$ & velocity field of the flow-matching denoiser \\
  \end{tabular}
\end{table}

\chapter{List of Abbreviations}

\begin{table}[h]
  \centering
  \begin{tabular}{ll}
    3DGS  & 3D Gaussian Splatting \\
    NeRF  & Neural Radiance Field \\
    VAE   & Variational Autoencoder \\
    ELBO  & Evidence Lower Bound \\
    KL    & Kullback--Leibler (divergence) \\
    SH    & Spherical Harmonics \\
    DiT   & Diffusion Transformer \\
    RoPE  & Rotary Position Embedding \\
    APE   & Absolute Position Embedding \\
    MLP   & Multilayer Perceptron \\
    FPS   & Farthest-Point Sampling \\
    InfoNCE & Information Noise-Contrastive Estimation \\
    PCA   & Principal Component Analysis \\
    SSC   & Semantic Scene Completion \\
    PSNR  & Peak Signal-to-Noise Ratio \\
    SSIM  & Structural Similarity Index \\
    LPIPS & Learned Perceptual Image Patch Similarity \\
    DC/AC & DC (mean) and AC (residual) components of colour \\
  \end{tabular}
\end{table}


\twocolumn
\mainmatter
        \chapter{Introduction}
\label{chap:intro}

%% ──────────────────────────────────────────────────────────────────────────
\section{Motivation}
\label{sec:intro_motivation}

The ability to generate or complete three-dimensional scenes is one of the
more ambitious goals in computer vision and graphics. Applications span
rapid content creation, embodied-AI simulation, autonomous navigation testing,
and robot planning in unseen environments. Progress towards this goal has
accelerated in recent years, partly because the representation of 3D scenes has
matured considerably.

3D Gaussian Splatting~\citep{kerbl2023gaussian} has become the dominant
representation for photorealistic scene reconstruction. A scene is modelled as
a set of anisotropic Gaussian splats, each carrying a position, an opacity, a
colour, and an anisotropic covariance parameterised by a rotation and a scale.
The representation renders in real time, matches or exceeds the quality of
earlier implicit approaches~\citep{mildenhall2020nerf}, and is explicitly
structured, meaning a network can in principle produce the primitives
directly rather than inferring them from a coordinate query.

Yet despite this progress in representation, scene-level 3DGS \emph{generation}
remains largely unexplored. The standard recipe for generation in neighbouring
domains is well-established: first compress the raw signal into a smooth,
compact latent space with a variational autoencoder~\citep{kingma2014vae},
then train a generative model on that latent~\citep{rombach2022ldm,
ho2020ddpm,lipman2023flowmatching}. This recipe has been proven at the object
level for 3D content~\citep{zeng2022lion,zhang2024gaussiancube,zhao2023michelangelo},
but scaling it to whole scenes is a different problem. A scene contains
hundreds of thousands of unordered Gaussian primitives, spans an unbounded
volume, mixes dozens of object categories, and comes with no canonical
coordinate frame or ordering. Building the scene-level VAE that would underpin
this recipe is the central challenge this thesis addresses.

Can3Tok~\citep{gao2025can3tok} established for the first time that a
scene-level 3DGS VAE is feasible: it demonstrated convergence on small training
sets and showed that the resulting latent could be used for generation.
This was an important existence proof, but it left open a question that became
the organising challenge of this thesis: can such an autoencoder be scaled
to thousands of heterogeneous scenes, and what is required to achieve that?

During the early stages of this project, the answer to that question proved
to be \emph{no, not without significant changes}. Training the Can3Tok
architecture on thousands of pre-fitted, frozen 3DGS scenes failed to
converge to anything useful. Understanding why, and fixing it, is what this
thesis is about.

%% ──────────────────────────────────────────────────────────────────────────
\section{Problem Statement}
\label{sec:intro_problem}

This thesis studies the construction and training of a semantically aware,
scalable, scene-level 3DGS variational autoencoder as Stage~1 of a two-stage
generative pipeline for scene generation and completion. The encoder must map
a variable-size, unordered Gaussian scene into a fixed-size structured latent
$Z \in \bR^{T \times d}$, and the decoder must reconstruct the scene from
it. Downstream, a flow-matching Stage~2 model should be trainable on the
resulting latent distribution.

The central difficulty we identify is that the per-primitive parameters of a
\emph{pre-fitted} 3DGS scene are structurally ambiguous in two compounding
ways.

\paragraph{The ordering problem.} A 3DGS scene is a set; any permutation of
its elements represents the same scene. A decoder that is asked to reconstruct
``the $i$-th Gaussian'' is being asked to reconstruct a target that has no
stable identity from one scene to the next. Without a canonical ordering,
the slot-aligned reconstruction loss has no persistent signal and the model
cannot converge.

\paragraph{The placement and orientation problem.} Even within a small spatial
region, the exact position of any individual Gaussian is largely a degree of
freedom of the fitting optimiser. And the covariance
$\vsigma_i = R(q_i)\,\mathrm{diag}(s_i^2)\,R(q_i)^\top$ is invariant to
sign flips and permutations of the quaternion and scale, meaning a naive loss
on the raw quaternion $q_i$ penalises differences that do not change the actual
Gaussian shape.

Both problems become much harder at scale. At 300 scenes the model has enough
repetition to memorise individual scenes despite these ambiguities. At 3,000
or 6,000 scenes, memorisation is no longer possible and the problems become
the dominant obstacle.

There is also a subtler interaction that we discovered empirically: the
\emph{sampling budget}. In Can3Tok, the Gaussian splats are jointly optimised
alongside the autoencoder, so the fitting process can co-adapt to the
tokenisation. With pre-fitted, frozen splats there is no such co-adaptation.
As a result, the number of Gaussians per latent token (the ``block size''
$g = \lceil N/T \rceil$) directly controls how much structural ambiguity the
decoder must resolve per token, and using a smaller budget is substantially
better in this regime.

%% ──────────────────────────────────────────────────────────────────────────
\section{Research Questions}
\label{sec:intro_rqs}

The work addresses one main question and three sub-questions.

\medskip
\noindent\textbf{Main question}: \textit{What is required to scale a
scene-level 3DGS variational autoencoder from a few hundred training scenes
to several thousand heterogeneous scenes, and which design choices matter most?}

\medskip
\noindent\textbf{RQ1 (Ordering and identifiability).} Does imposing a canonical
spatial order on the Gaussian set restore convergence at scale, and how do
different ordering strategies compare?

\noindent\textbf{RQ2 (Architecture and generalisation).} Does replacing the
global encoder bottleneck with a spatially local encoder close the
train-to-test generalisation gap that appears at scale?

\noindent\textbf{RQ3 (Semantic structure).} Can per-primitive semantic
supervision and latent disentanglement be incorporated into the small-scale
regime without degrading reconstruction, motivating their inclusion in the
final design?

%% ──────────────────────────────────────────────────────────────────────────
\section{Contributions}
\label{sec:intro_contributions}

The contributions of this thesis are as follows.

\begin{enumerate}
\item \textbf{Identification of the scaling barriers.} We show, through
  systematic experimentation, that the ordering of the Gaussian set is the
  primary obstacle to scaling. Without a stable slot-to-position correspondence,
  the reconstruction target is not learnable at thousands of scenes.

\item \textbf{Fixed-frame Hilbert serialisation.} We apply a space-filling
  Hilbert curve in a fixed canonical coordinate frame to impose a consistent,
  locality-preserving order on the Gaussian set~\citep{skilling2004hilbert,
  wu2024ptv3}. We show that this is more effective than opacity ranking
  (the Can3Tok default) and Z-order (Morton) sorting.

\item \textbf{Anchor-relative position decoding.} We introduce per-token
  spatial anchors (block centroids) and decode each Gaussian's position as
  $\hat{\vmu}_i = a_{\tau(i)} + \rho\,\tanh(o_i)$, so the decoder predicts
  a bounded local offset rather than an absolute coordinate. This matches the
  Hilbert block structure and substantially reduces the within-block structural
  ambiguity the decoder must resolve.

\item \textbf{Improved reconstruction objectives.} We combine a Sinkhorn
  entropic-OT set loss~\citep{cuturi2013sinkhorn} that matches primitives
  within blocks by optimised assignment, with a Bures co-diagonalisation
  covariance loss that measures orientation and shape in a manner invariant to
  quaternion parameterisation.

\item \textbf{Spatially local encoder and token-local decoder.} We replace the
  globally-mixed bottleneck with a windowed cross-attention encoder and a
  shared per-token MLP decoder. We show that this is the single change most
  responsible for closing the train-to-test generalisation gap.

\item \textbf{Small-scale semantic exploration.} We conduct 300-scene
  experiments exploring per-primitive contrastive objectives and latent
  disentanglement, demonstrating that these are feasible and produce
  meaningful semantic structure in the latent.

\item \textbf{Stage-2 generative design.} We specify and partially implement
  a hierarchical flow-matching generator~\citep{lipman2023flowmatching,
  peebles2023dit} built on the Stage-1 latent, targeting both unconditional
  scene generation and partial-scan completion.
\end{enumerate}

We are upfront about scope: the large-scale ablation training is underway
at the time of writing and full quantitative results will be incorporated as
they become available. The Stage-2 training has not yet completed, and scene
completion remains the primary direction for future work.

%% ──────────────────────────────────────────────────────────────────────────
\section{Thesis Structure}
\label{sec:intro_structure}

Chapter~\ref{chap:related} positions this work in the literature.
Chapter~\ref{chap:method} presents the full method in detail.
Chapter~\ref{chap:results} presents the experimental findings: the 300-scene
exploratory results first, then the large-scale comparisons with
placeholders for in-progress runs. Chapter~\ref{chap:discussion} interprets
the findings and discusses limitations. Chapter~\ref{chap:conclusion}
answers the research questions. The appendix records all configuration details
for reproducibility.
        \chapter{Related Work}
\label{chap:related}

This chapter situates the thesis in four bodies of work. I cover explicit 3D scene representations and where 3DGS sits among them; latent generative modelling of 3D content, with attention to the two-stage VAE+generator recipe; learning on unordered sets and the use of space-filling curves to impose a canonical order; and contrastive plus disentangled representation learning, which informs my small-scale semantic experiments. Each section closes by stating the specific gap the thesis addresses.

%% ──────────────────────────────────────────────────────────────────────────
\section{3D Scene Representations}
\label{sec:rw_repr}

\paragraph{Neural radiance fields.}
NeRF \citep{mildenhall2020nerf} represents a scene as a continuous volumetric function encoded in the weights of a coordinate-conditioned MLP. It produces photorealistic novel views and remained the dominant approach for several years. Two limitations hold it back as a generative target: rendering is slow because each ray needs hundreds of network evaluations, and the per-scene weight representation has no explicit parameter for a network to predict on a scene level.

\paragraph{3D Gaussian Splatting.}
3DGS \citep{kerbl2023gaussian} represents a scene as a set of anisotropic Gaussian primitives rasterised by alpha compositing. The approach renders in real time, matches NeRF on novel-view synthesis, and exposes an explicit per-primitive parameterisation. Each primitive has a position $\mu \in \bR^3$, a colour $c \in \bR^3$ (the DC spherical-harmonic coefficient, ignoring higher-order view-dependent terms here), an opacity $\alpha \in (0,1)$, a per-axis scale $s \in \bR^3_+$, and a unit quaternion $q$ defining the orientation. The covariance matrix is then $\Sigma = R(q)\,\mathrm{diag}(s^2)\,R(q)^\top$.

Scaffold-GS \citep{lu2024scaffold} reduces redundancy in a dense 3DGS fit by anchoring per-primitive attributes to a sparser neural scaffold, predicting offsets relative to anchor points rather than absolute positions. This anchor-plus-offset structure is what motivates the position decoding choice in \S\ref{sec:method_anchor}. GaussianCube \citep{zhang2024gaussiancube} takes a different route: it fits a fixed Gaussian count at the object level and then transports the primitives onto a regular voxel grid, so a standard convolutional diffusion model can operate on a 3D tensor. Both of these object-level approaches inform my own per-token anchor decoding, but neither extends directly to unbounded indoor scenes.
\paragraph{Gap.} 3DGS gives an explicit representation, but turning it into a generative target at scene level requires a compact, organised latent. Neither 3DGS itself nor Scaffold-GS provides that latent; building it is the goal of Stage~1 in this thesis.

%% ──────────────────────────────────────────────────────────────────────────
\section{Latent Generative Modelling of 3D Content}
\label{sec:rw_gen}

\paragraph{The two-stage recipe.}
A widely used recipe for high-resolution generation is to first train a VAE to compress the signal into a latent \citep{kingma2014vae}, and then train a generative model in that latent space \citep{rombach2022ldm}. The generative half has shifted from DDPM \citep{ho2020ddpm} to flow matching \citep{lipman2023flowmatching, liu2023rectified}, both typically realised with a transformer backbone such as DiT \citep{peebles2023dit} or SiT \citep{ma2024sit}. The latent decouples the awkward structure of the raw signal from the easier learning problem of modelling a tame distribution of token sequences.

\paragraph{Object-level 3D generation.}
LION \citep{zeng2022lion} applies hierarchical latent diffusion to point cloud objects. Michelangelo \citep{zhao2023michelangelo} uses a cross-attention Perceiver \citep{jaegle2021perceiver, jaegle2022perceiverio} to encode a 3D shape into a fixed set of latent tokens aligned with image and text embeddings; its shape-latent VAE is the architectural ancestor of mine. GaussianCube \citep{zhang2024gaussiancube} structures the per-primitive prediction into a position-plus-offset form so a dense convolutional model can operate on it. All three are object-level, with bounded extent and small fixed primitive counts.

\paragraph{Scene-level 3DGS.}
Can3Tok \citep{gao2025can3tok} is the direct predecessor of this thesis. It uses a Perceiver-style cross-attention encoder to compress a scene's Gaussian primitives into a fixed-size latent, and a transformer-plus-MLP decoder to reconstruct the per-primitive parameters. Can3Tok was demonstrated on a few hundred scenes and showed conditional generation as a downstream application. SceneSplat \citep{li2025scenesplat} introduces the first generalisable feed-forward 3DGS encoder, pre-trained on its accompanying SceneSplat-7K dataset of 7{,}916 indoor 3DGS scenes; the dataset comes with per-Gaussian ScanNet-72 \citep{dai2017scannet} category labels. SceneSplat-7K is the training corpus throughout this thesis. SceneSplat++ \citep{ma2025scenesplatpp} extends the dataset and benchmarks language-aligned Gaussian Splatting more thoroughly.

\paragraph{Gap.} What is required to scale a scene-level 3DGS VAE from hundreds to several thousand pre-fitted scenes, and which architectural choices control whether it converges, has not been systematically studied. The seven-experiment ablation in Chapter~\ref{chap:results} is my answer.

%% ──────────────────────────────────────────────────────────────────────────
\section{Learning on Unordered Sets and Space-Filling Curves}
\label{sec:rw_sets}

\paragraph{Permutation-invariant encoders.}
PointNet \citep{qi2017pointnet} established symmetric aggregation as a way to produce permutation-invariant features over point sets. The Perceiver \citep{jaegle2021perceiver} generalises this by cross-attending a small learned query set against the input, yielding a fixed-size latent that is invariant to input permutation. Permutation invariance is cheap on the encoder side; the difficulty reappears on the decoder side, where any element-wise reconstruction loss compares predicted slot $i$ to target slot $i$ and so implicitly fixes an ordering.

\paragraph{Canonical orderings via space-filling curves.}
Rather than handle permutation in the loss, a complementary approach is to fix a canonical order on the data. Point Transformer V3 \citep{wu2024ptv3} serialises point cloud tokens along the Hilbert space-filling curve \citep{skilling2004hilbert} and shows that this ordering, with its tight locality property (consecutive indices remain spatially adjacent), gives better downstream results than the alternatives such as Z-order (Morton) curves, where high-order bit flips cause long spatial jumps. The same idea, but pushed onto the reconstruction \emph{target}, is the central mechanism behind RQ1 in this thesis.

\paragraph{Gap.} Canonical serialisation is well known on the processing side of 3D learning but, to the best of my knowledge, has not been used with a fixed-frame coordinate system to make the per-slot 3DGS regression target identifiable for a generative VAE. Doing this in concert with a per-token spatial anchor is what makes the slot-aligned loss learnable in my setup.

%% ──────────────────────────────────────────────────────────────────────────
\section{Semantic and Disentangled Representations}
\label{sec:rw_sem}

\paragraph{Contrastive learning.}
Self-supervised contrastive learning \citep{oord2018cpc, chen2020simclr} learns features by pulling together related samples (positives) and pushing apart unrelated ones (negatives) in embedding space. The InfoNCE objective \citep{oord2018cpc} is the standard tool. I adapt it in two forms in this thesis: a per-Gaussian variant that uses the per-Gaussian ScanNet-72 labels from SceneSplat-7K as the positive criterion, and a scene-level variant that uses the scene's label-distribution histogram.

\paragraph{Disentangled latent representations.}
$\beta$-VAE \citep{higgins2017betavae} attempts to disentangle independent factors of variation in the latent by reweighting the KL term in the VAE objective. Locatello et al.\ \citep{locatello2019challenging} show that unsupervised disentanglement is in general unidentifiable without inductive bias, which is why my small-scale experiments rely on structural inductive bias (a layout/geometry split with auxiliary heads and cross-reconstruction) rather than on KL reweighting alone.

\paragraph{Gap.} Per-Gaussian semantic labels for 3DGS scenes have become available with SceneSplat-7K \citep{li2025scenesplat}. So far they have been used for discriminative scene understanding rather than to add semantic structure to a generative latent. My 300-scene exploration in \S\ref{sec:results_300} is a first look at how that supervision interacts with a Stage~1 latent designed for downstream generation.
        \chapter{Methodology}
\label{sec:methodology}

This chapter presents the method. We first formalize the problem and the notion of identifiability that organizes the thesis (\autoref{sec:method_overview}). We then describe the shared encoder (\autoref{sec:method_encoder}), the family of latent-structure variants (\autoref{sec:method_latent}), the spatially local encoder that is our main fix for scaling (\autoref{sec:method_local}), and the family of decoders (\autoref{sec:method_decoder}). We then give the reconstruction objectives, including the gauge-invariant covariance and permutation-invariant set losses (\autoref{sec:method_recon}), and the semantic and disentanglement objectives (\autoref{sec:method_semantic}). Finally we describe the data pipeline (\autoref{sec:method_data}), the second-stage generative design (\autoref{sec:method_stage2}), and the experimental setup (\autoref{sec:method_setup}). Throughout, we are explicit about which components are inherited from Can3Tok \citep{gao2025can3tok} and Michelangelo \citep{zhao2023michelangelo} and which are our additions.

% =====================================================================
\section{Overview and Problem Formulation}
\label{sec:method_overview}
% =====================================================================

\paragraph{The representation.} A scene is a set of $N$ 3D Gaussian primitives $\mathcal{G} = \{g_i\}_{i=1}^{N}$. Each primitive is
\begin{equation}
  g_i = (\mu_i,\, c_i,\, \alpha_i,\, s_i,\, q_i),
\end{equation}
with position $\mu_i \in \mathbb{R}^3$, colour $c_i \in \mathbb{R}^3$ (the diffuse term of the spherical-harmonic colour), opacity $\alpha_i \in (0,1)$, per-axis scale $s_i \in \mathbb{R}^3_{+}$, and unit quaternion $q_i \in \mathbb{S}^3$. The orientation and scale define the anisotropic covariance
\begin{equation}
  \Sigma_i \;=\; R(q_i)\,\mathrm{diag}(s_i^2)\,R(q_i)^\top,
  \label{eq:covariance}
\end{equation}
where $R(q)$ is the rotation matrix of $q$. We pack each primitive into a $14$-dimensional vector and a scene into a tensor in $\mathbb{R}^{N \times 14}$.

\paragraph{The goal.} We want an autoencoder, the first stage of a two-stage pipeline, consisting of an encoder $E_\phi$ that maps $\mathcal{G}$ to a fixed-size latent $Z \in \mathbb{R}^{T \times d}$ ($T$ tokens of width $d$) and a decoder $D_\theta$ that reconstructs $\hat{\mathcal{G}}$ from $Z$. We train it as a variational autoencoder \citep{kingma2014vae}: $E_\phi$ outputs the parameters of a diagonal Gaussian posterior $q_\phi(Z \mid \mathcal{G}) = \mathcal{N}(\mu_Z, \sigma_Z^2)$, we sample $Z$ by the reparameterization trick, and we optimize a reconstruction term against a Kullback-Leibler regularizer toward a standard normal prior,
\begin{equation}
  \mathcal{L} \;=\; \mathcal{L}_{\mathrm{recon}}(\hat{\mathcal{G}}, \mathcal{G}) \;+\; \beta\, D_{\mathrm{KL}}\!\big(q_\phi(Z\mid\mathcal{G}) \,\|\, \mathcal{N}(0, I)\big),
  \label{eq:elbo}
\end{equation}
with a deliberately small $\beta$ so that the latent retains enough capacity for accurate reconstruction while remaining smooth enough for the second stage. The second stage then learns a generative model of $Z$ (\autoref{sec:method_stage2}). A schematic of the full pipeline is given in \autoref{fig:pipeline}.

\begin{figure*}[t]
  \centering
  \includegraphics[width=0.9\linewidth, height=0.26\textheight]{example-image}
  \caption[Two-stage pipeline overview]{\textbf{Two-stage pipeline.} Stage 1 (this thesis) is a variational autoencoder that tokenizes a scene-level 3DGS set into a structured latent $Z$ and reconstructs the Gaussians. Stage 2 is a latent flow-matching transformer that models $p(Z)$ for generation, and an inpainting variant for scene completion. \emph{[FIGURE TODO: replace with the exported pipeline diagram. Show: raw 3DGS scene $\to$ normalization and Hilbert serialization $\to$ shared encoder $\to$ latent $Z$ $\to$ decoder $\to$ reconstructed Gaussians, with the Stage 2 layout/geometry models feeding the frozen decoder. Source: README ``Architecture at a Glance'' and ``Second-Stage Generation Pipeline''.]}}
  \label{fig:pipeline}
\end{figure*}

\paragraph{Identifiability and gauge.} The central difficulty is that the per-primitive parameters of a \emph{pre-fitted} 3DGS scene are non-identifiable: distinct parameter sets render to the same image and are therefore equivalent as scene content, yet differ as regression targets. We call this nuisance variation \emph{gauge}. Three facets recur.
\begin{itemize}
  \item \textbf{Set gauge (order).} $\mathcal{G}$ is a set; any permutation of its elements is the same scene. A loss that compares the $i$-th output to the $i$-th target therefore depends on an arbitrary ordering and is not a well-defined function of the scene.
  \item \textbf{Orientation gauge.} The covariance in \autoref{eq:covariance} is invariant to the quaternion double cover ($q \equiv -q$), to relabelling the axes (permuting the columns of $R$ together with the entries of $s$), to column sign flips, and, when $s$ is near-isotropic, to almost any rotation. An element-wise loss on $q$ penalizes all of this, none of which changes the Gaussian.
  \item \textbf{Placement gauge.} Within a small spatial region, exactly where each individual Gaussian sits is largely a freedom of the fit; regressing the exact position forces the network to memorize nuisance.
\end{itemize}
Our thesis is that scaling fails primarily because these gauges make the naive target unlearnable, and that the fixes are to (i) impose a canonical order, (ii) replace gauge-sensitive losses with gauge-invariant ones, and (iii) make the latent compositional so it generalizes rather than memorizes. The remainder of the chapter develops each.

% =====================================================================
\section{Shared Encoder}
\label{sec:method_encoder}
% =====================================================================

All variants share one encoder, inherited in form from the Michelangelo shape-latent VAE \citep{zhao2023michelangelo} and Can3Tok \citep{gao2025can3tok}, and built on a Perceiver-style cross-attention bottleneck \citep{jaegle2021perceiver}.

\paragraph{Input features.} Each Gaussian is lifted to an input feature by concatenating a Fourier positional encoding \citep{tancik2020fourier} of its normalized absolute position, a Fourier encoding of a coarse voxel-centre coordinate (a hashed grid index that gives the encoder an explicit notion of cell membership), and the raw Gaussian parameters (opacity, scale, quaternion, and colour or colour residual). A linear projection maps the concatenated feature to the model width.

\paragraph{Cross-attention bottleneck.} A set of learned query tokens cross-attends to the projected Gaussian features. We use one global query, whose output we call the \emph{shape embedding} $e \in \mathbb{R}^{w}$ and which summarizes the whole scene, and $K$ geometry queries, whose outputs are the per-region geometry tokens. In the baseline encoder a stack of self-attention layers then mixes the query tokens. Because the queries are fixed in number, the latent size is independent of $N$: a scene with any number of primitives is encoded into the same token budget, which is the property that lets the encoder ingest a variable-size, unordered set \citep{jaegle2022perceiverio}. The encoder is therefore permutation-invariant by construction, which already resolves the set gauge on the \emph{input} side; the set gauge re-appears only in the decoder objective (\autoref{sec:method_recon}).

% =====================================================================
\section{Latent-Structure Variants}
\label{sec:method_latent}
% =====================================================================

We study three ways of organizing the latent $Z$, illustrated in \autoref{fig:latent_variants}. They differ in whether the latent is undifferentiated, explicitly split into semantic and geometric subspaces, or kept in one-to-one correspondence with spatial regions.

\begin{figure*}[t]
  \centering
  \includegraphics[width=0.95\linewidth, height=0.28\textheight]{example-image}
  \caption[Latent-structure and decoder variants]{\textbf{Stage 1 architecture variants.} Left: the baseline (C) with undifferentiated geometry tokens. Centre: the disentangled latent (A) splitting $Z$ into a semantic subspace $z_s$ and a geometry subspace $z_g$ with auxiliary heads and disentanglement losses. Right: the structured per-token latent with the spatially local encoder, where token $k$ corresponds to Hilbert block $k$. Decoder options (flat MLP, token-local MLP, anchor-relative, cross-attention conditioning) are shown beneath. \emph{[FIGURE TODO: replace with exported diagrams from README sections ``Shared Encoder Pipeline'', ``Strategy A/B1/C'', and the structured/local-encoder description. Keep the three latent layouts side by side.]}}
  \label{fig:latent_variants}
\end{figure*}

\paragraph{Baseline (undifferentiated).} The simplest variant treats $Z$ as $T$ undifferentiated geometry tokens with no semantic split and no layout conditioning. This is the closest analogue to the published Can3Tok latent and serves as our control.

\paragraph{Disentangled latent.} Motivated by the wish for a latent in which coarse scene identity and fine geometry are separated, we split $Z$ into a small semantic subspace $z_s$ (a handful of tokens routed from the shape embedding) and a larger geometry subspace $z_g$ (routed from the geometry tokens). Both are KL-regularized. Auxiliary heads read $z_s$ to predict scene-level quantities (a mean colour, a category-distribution, and per-category spatial centroids), which encourages $z_s$ to carry layout-level information, while disentanglement losses (\autoref{sec:method_semantic}) push $z_g$ to be sufficient for geometry and to be statistically separated from $z_s$. This variant is where our semantic-structure and disentanglement experiments were carried out at small scale.

\paragraph{Structured per-token latent.} The third variant, which is the one used in our large-scale model, keeps the latent in one-to-one correspondence with space. After serializing the Gaussian set along a Hilbert curve (\autoref{sec:method_data}), a contiguous block of primitives is a local spatial cluster. Latent token $k$ is made to encode block $k$: the per-token posterior is used directly as $z$, and the dense global re-projection that the baseline uses to mix all tokens is skipped. This is what makes the latent \emph{compositional}: token $k$ always describes the same kind of local patch, patches recur across scenes, and so the representation transfers rather than memorizing a per-scene global code. Realizing this fully requires the local encoder of \autoref{sec:method_local}.

% =====================================================================
\section{Spatially Local Encoder}
\label{sec:method_local}
% =====================================================================

The baseline encoder mixes all query tokens with global self-attention, so each latent token is a function of the whole scene. At small scale this is harmless, but at the scale of thousands of heterogeneous scenes we found it encourages the latent to behave as a memorized global code, which generalizes poorly (\autoref{sec:results}). We therefore replace the global bottleneck with a windowed, spatially local encoder.

Concretely, with the Gaussian set Hilbert-ordered into $K$ contiguous blocks of $g = \lceil N/K \rceil$ primitives each, geometry query $k$ attends only to the primitives in a local window of blocks around block $k$ (from block $k-\omega$ to block $k+\omega$ for a small half-width $\omega$), rather than to the whole scene. A single global query still attends to all primitives to provide a scene-level summary. No global self-attention runs inside the encoder, since that would re-mix the tokens and destroy the locality the windowing creates; cross-token interaction is deferred to the decoder transformer. Latent token $k$ therefore encodes the local geometry of block $k$, the same block that the decoder's anchor and token-local head reconstruct, giving an end-to-end alignment between encoder token, latent token, spatial region, and decoder slot. We treat the window scope as an ablation axis: setting the window to span the whole scene recovers a global encoder with otherwise identical structure, which isolates locality as the variable of interest (\autoref{sec:results}).

% =====================================================================
\section{Decoders}
\label{sec:method_decoder}
% =====================================================================

The decoder maps $Z$ back to $N$ Gaussians. We consider several options that trade capacity against structure.

\paragraph{Flat decoder.} The inherited decoder applies a transformer to the post-projected latent and then a large multilayer perceptron that flattens all tokens and regresses all $N \times 14$ outputs at once. This has very high capacity but a single dense bottleneck, which we found encourages memorization at scale.

\paragraph{Token-local decoder.} We replace the flat head with a shared per-token multilayer perceptron: each of the $T$ decoder tokens decodes its own contiguous slice of $g$ Gaussians using shared weights. This removes the dense bottleneck, is far smaller in parameter count, and matches the block structure of the latent so that token $k$ decodes the primitives of block $k$. Per-primitive semantic features for the contrastive objective are routed through a pooled-hidden path that preserves the interface the flat decoder exposed.

\paragraph{Anchor-relative position decoding.} In the spirit of Scaffold-GS \citep{lu2024scaffold} and the structural-reference idea of GaussianCube \citep{zhang2024gaussiancube}, we decode position as an anchor plus a bounded local offset,
\begin{equation}
  \hat{\mu}_i \;=\; a_{\,\tau(i)} \;+\; \rho \cdot \tanh(o_i),
  \label{eq:anchor}
\end{equation}
where $a_{\tau(i)}$ is the anchor of the token $\tau(i)$ that owns primitive $i$ (the per-block centroid, predicted from the token), $o_i$ is the decoder's raw position output, and $\rho > 0$ is a learnable global offset scale. The $\tanh$ bound is the point: it forces the offset to be a genuinely \emph{local} displacement, so the coarse position must be carried by the anchor. This directly targets the placement gauge: the network is asked for ``which block, plus a small local correction'' rather than an absolute coordinate to be memorized.

\paragraph{Conditioning decoders.} For the disentangled latent we also study decoders in which a layout subspace conditions geometry. In the cross-attention variant, the geometry tokens form the decoder sequence and the layout tokens are injected as keys and values at every layer, which is a stronger architectural separation than letting all tokens mix by self-attention. A deterministic projection of the shape embedding can also serve as a stable layout code for conditioning, which is useful for completion because encoding the same partial scene twice yields the same conditioning. These variants are most relevant to the second-stage design (\autoref{sec:method_stage2}).

% =====================================================================
\section{Reconstruction Objectives}
\label{sec:method_recon}
% =====================================================================

The reconstruction term in \autoref{eq:elbo} is where the gauge problem is won or lost. We describe the objectives in increasing order of gauge-awareness.

\paragraph{Element-wise loss with serialized order.} The simplest reconstruction loss compares matched slots,
\begin{equation}
  \mathcal{L}_{\mathrm{elem}} \;=\; \frac{1}{B}\sum_{i=1}^{N} \big\| \hat{g}_i - g_i \big\|_2,
\end{equation}
optionally with per-attribute weights. As discussed, this is only a meaningful function of the scene if slot $i$ corresponds to a stable target. We obtain such a correspondence by serializing both prediction and target along a fixed-frame Hilbert curve (\autoref{sec:method_data}), so that slot $i$ is a spatially stable location rather than, for example, the $i$-th most opaque primitive. This converts the otherwise unlearnable slot-to-target map into a Lipschitz one, the property for which the Hilbert curve is preferred over the Z-order curve \citep{wu2024ptv3}.

\paragraph{Permutation-invariant set loss.} Serialization fixes the order \emph{between} blocks but not \emph{within} a block, where the assignment of which colour or orientation belongs to which primitive remains a gauge. A slot-aligned loss forces the decoder to break that gauge, which it cannot, so it emits the per-block mean (washed-out colour, round blobs). We therefore optionally match predicted and target primitives \emph{within} each block by entropic optimal transport. Writing $C$ for the matrix of squared feature distances between the $g$ predicted and $g$ target primitives of a block, we compute a soft assignment by Sinkhorn iteration \citep{cuturi2013sinkhorn} and score the matched pairs, following the assignment-then-score recipe of DETR \citep{carion2020detr}. The matching is detached, so the gradient flows only through the matched residual. The identity assignment recovers $\mathcal{L}_{\mathrm{elem}}$ exactly, so the set loss can only lower it.

\paragraph{Gauge-invariant covariance loss.} For orientation and scale we avoid an element-wise loss on the quaternion, which penalizes the orientation gauge of \autoref{sec:method_overview}, and instead compare the covariances in \autoref{eq:covariance} directly. Two covariances that differ only by gauge are equal as matrices, so any metric on the matrix is gauge-invariant. We use the Bures (equivalently, $2$-Wasserstein between zero-mean Gaussians) shape term \citep{villani2009ot},
\begin{equation}
  d_{\mathrm{B}}^2(\Sigma_p, \Sigma_t) \;=\; \mathrm{tr}\!\Big( \Sigma_p + \Sigma_t - 2\,\big( \Sigma_t^{1/2}\,\Sigma_p\,\Sigma_t^{1/2} \big)^{1/2} \Big),
  \label{eq:bures}
\end{equation}
and, as a cheaper and numerically very stable surrogate, the co-diagonalization upper bound
\begin{equation}
  d_{\mathrm{cd}}^2(\Sigma_p, \Sigma_t) \;=\; \big\| \Sigma_p^{1/2} - \Sigma_t^{1/2} \big\|_F^2,
  \label{eq:burescodiag}
\end{equation}
which equals \autoref{eq:bures} when the two matrices commute. Matrix square roots are computed by a Newton-Schulz iteration under trace normalization, which is pure matrix multiplication and therefore differentiable and stable even at coincident eigenvalues. Position, colour, and opacity keep their ordinary residual norms, so only the scale-and-rotation block is replaced and the overall loss scale (and hence the balance with the KL and auxiliary terms) is preserved. Because the covariance square root is available in closed form as $\Sigma^{1/2} = R\,\mathrm{diag}(|s|)\,R^\top$, the gauge-invariant shape term folds directly into the per-primitive feature used by the Sinkhorn matching, so the set loss and the covariance loss compose at no extra cost.

\paragraph{Diagnostics.} Because the raw quaternion error is gauge-blind, we monitor a gauge-invariant covariance error (mean per-primitive Bures distance) and the target anisotropy at evaluation time, so that reconstruction quality of orientation is measured by a quantity that is actually well-defined.

% =====================================================================
\section{Semantic and Disentanglement Objectives}
\label{sec:method_semantic}
% =====================================================================

To make the latent semantically structured we add objectives that consume the per-primitive labels of SceneSplat-7K \citep{li2025scenesplat}.

\paragraph{Per-primitive contrastive loss.} We attach a projection head that produces a normalized feature for each Gaussian and apply an InfoNCE objective \citep{oord2018cpc,chen2020simclr} that pulls together primitives sharing a category label and pushes apart those that do not, implemented efficiently against per-category prototypes. This is the objective that organizes the per-primitive feature space; we visualize its effect by principal-component projection of the features (\autoref{sec:results}). Only the labelled scenes contribute to this term; scenes from auxiliary sources without compatible labels are filtered out so there is no cross-taxonomy collision.

\paragraph{Scene-level and disentanglement objectives.} For the disentangled latent we additionally use: auxiliary heads on $z_s$ that predict scene-level quantities; a cross-reconstruction loss that decodes a latent with $z_s$ swapped across the batch but $z_g$ kept, and requires the geometry to be reconstructed correctly, which forces $z_g$ to be sufficient for geometry independently of $z_s$; and an orthogonality penalty that decorrelates random projections of the $z_s$ and $z_g$ means. Consistent with the identifiability literature \citep{locatello2019challenging}, we treat this as weak, supervised structuring of the latent rather than strict factorization.

% =====================================================================
\section{Data Pipeline}
\label{sec:method_data}
% =====================================================================

The data pipeline is where the gauge fixes are partly realized, and where the sampling question that we identify as central is decided.

\paragraph{Dataset.} We use SceneSplat-7K \citep{li2025scenesplat}, comprising $7{,}916$ indoor 3DGS scenes derived from seven established sources, with per-primitive ScanNet-style category labels. Our large-scale training mixes a set of chunked interior scenes (which carry semantics) with additional full-scene sources (whose semantics are disabled and which therefore contribute only to reconstruction), and validates on held-out full scenes and on held-out chunks; the latter, being disjoint by construction from the training chunks, gives a clean in-distribution generalization measure. Exact counts and paths are given in \autoref{sec:method_setup} and the appendix.

\paragraph{Normalization.} Scenes are unbounded and scale-inconsistent, so each is normalized into a canonical sphere of fixed radius. For chunked scenes we compute the normalization from the union of all chunks of a scene, so that all chunks of a room share one global coordinate frame, which is what makes a position target consistent across chunks of the same scene; full scenes use a per-scene normalization into the same radius.

\paragraph{Sampling: the central practical lever.} A scene has far more primitives than we can encode, so a fixed number $N$ must be sampled. We study two choices. Opacity sampling keeps the $N$ most opaque primitives; this is simple and visually clean but makes the sampled density follow the opacity field, which is non-uniform and hard to compress. Density-uniform (farthest-point-style) sampling instead keeps a spatially uniform subset, which equalizes the per-token load. We also vary $N$ itself. This matters because the original Can3Tok pipeline samples a large budget (on the order of $40$k primitives) \emph{while the splats are still being jointly optimized}, so the optimization can adapt the primitives to the sample; in our setting the splats are pre-fitted and frozen, so the sample is of a fixed, gauge-ridden set. We found that a smaller budget (on the order of $10$k) converges substantially better than the larger one in the frozen-splat setting, which we attribute to the reduced within-block placement gauge per token. This sampling-and-ordering interaction is examined in \autoref{sec:results}.

\paragraph{Serialization.} After sampling, the primitives are reordered along a Hilbert space-filling curve computed in the fixed canonical frame, so that array slot $i$ is a spatially stable location and consecutive slots are spatial neighbours \citep{skilling2004hilbert,wu2024ptv3}. The same ordering defines the contiguous blocks used by the local encoder, the anchors, and the token-local decoder. We retain the Z-order (Morton) curve as an ablation to test the locality claim.

% =====================================================================
\section{Second-Stage Generative Model}
\label{sec:method_stage2}
% =====================================================================

The second stage learns to generate the Stage 1 latent. We adopt flow matching \citep{lipman2023flowmatching,liu2023rectified} with a transformer backbone \citep{peebles2023dit} operating on the frozen Stage 1 latent, and decode generated latents with the frozen Stage 1 decoder. The design is hierarchical, mirroring the latent's layout-versus-geometry structure, and is shown in \autoref{fig:stage2}.

\begin{figure*}[t]
  \centering
  \includegraphics[width=0.95\linewidth, height=0.27\textheight]{example-image}
  \caption[Stage 2 generation and completion variants]{\textbf{Stage 2 variants.} A layout model generates the coarse (layout/semantic) latent tokens from noise (optionally conditioned on text), then a geometry model generates the geometry tokens conditioned on the layout tokens; the frozen Stage 1 decoder turns the full latent into Gaussians. Scene completion replaces generation of the observed tokens with an inpainting procedure that fixes the encoded observed tokens and denoises only the unobserved ones. \emph{[FIGURE TODO: replace with exported diagrams from README ``Second-Stage Generation Pipeline'' (the layout DiT $\to$ geometry DiT generation flow and the completion/inpainting flow). Keep generation and completion as two rows.]}}
  \label{fig:stage2}
\end{figure*}

\paragraph{Generation.} A layout model is trained by flow matching to transport noise to the coarse latent tokens, optionally conditioned on a text or class token by cross-attention. A geometry model is then trained to transport noise to the geometry tokens conditioned on the coarse tokens. Conditioning the geometry model on the coarse tokens by the same cross-attention structure used in the Stage 1 conditioning decoder keeps the two stages aligned in the same representational space. At inference, the layout model samples coarse tokens, the geometry model samples geometry tokens conditioned on them, and the frozen decoder produces the scene, with no encoder and no ground truth required.

\paragraph{Completion.} Scene completion is posed as latent inpainting. A partial observation is encoded; the observed latent tokens are held fixed while the unobserved tokens are initialised from noise and denoised by the completion model, conditioned on the (stable) coarse tokens; the frozen decoder then produces the completed scene. A deterministic layout code is advantageous here because the same partial observation yields the same conditioning. We implement and train these models on the selected representations (\autoref{sec:res_selection}); in our setup the partial observation is simulated by masking a subset of the latent tokens of a fully encoded scene rather than by encoding a genuinely partial scan, which keeps the data pipeline fixed but means we measure latent-space inpainting rather than completion from a raw partial capture, a distinction we are explicit about in \autoref{sec:discussion}. Because each token corresponds to a Hilbert block, masking tokens hides spatial regions of the scene.

\paragraph{Two latent schemas.} The second stage operates on whichever latent the chosen Stage~1 model produces, and there are two cases. In the \emph{disentangled} schema the latent already separates into a small layout subspace $z_s$ and a larger geometry subspace $z_g$, so generation is the layout-then-geometry pair above and completion conditions the geometry on $z_s$. In the \emph{structured-flat} schema there is a per-token geometry latent but no $z_s$, so there is no layout model and completion runs without conditioning, denoising the masked tokens from the unmasked ones alone. Running completion in both schemas is what lets us isolate the contribution of the layout conditioning.

\paragraph{Positional encoding.} The flow-matching transformers need a notion of token position. We use a one-dimensional rotary embedding \citep{su2021roformer} for the geometry and completion models, whose tokens lie along the Hilbert curve and therefore carry a meaningful one-dimensional order, and a learned absolute embedding for the layout model, whose tokens are abstract and have no spatial order. This choice is supported by a small ablation in \autoref{sec:res_stage2}.

% =====================================================================
\section{Experimental Setup}
\label{sec:method_setup}
% =====================================================================

\paragraph{Architecture and training.} The encoder and decoder use a transformer width of $384$ with $12$ attention heads, $6$ encoder and $12$ decoder layers, and an $8$-frequency Fourier embedding, following the inherited configuration. The latent uses $512$ tokens of width $32$; the per-token width is the principal capacity knob. We train with AdamW, a base learning rate of $1\times 10^{-4}$ with cosine warm restarts and a short linear warmup, a small fixed KL weight ($10^{-5}$), mixed bf16 precision, and distributed data parallelism across H100 GPUs. The large-scale training set mixes $3{,}800$ semantically labelled interior chunks with three additional full-scene sources (roughly $800$, $1{,}290$, and $856$ scenes, semantics disabled), for about $6{,}700$ scenes in total; the Stage~2 models reuse this same combination so that the frozen encoder sees the distribution it was trained on. Full hyperparameters, the exact loss weights, the per-source counts, and the run commands are listed in \autoref{sec:apx:first_appendix} and released with the code for reproducibility.

\paragraph{Evaluation.} We evaluate reconstruction by per-attribute residual norms (position, colour, opacity, scale) reported on the same raw scale for training and validation so they are directly comparable, by the gauge-invariant covariance error for orientation and shape, and, for the final model, by image-space fidelity of renders of the reconstructed scenes (PSNR, SSIM, and LPIPS \citep{wang2004ssim,zhang2018lpips}); space is left for these in \autoref{sec:results} as the final model finishes. We validate on two held-out sets: a primary set of held-out full scenes (the harder, more representative case) and a disjoint set of held-out chunks, and we report both. As noted in \autoref{sec:res_scaling}, the full-to-chunk ratio is a coarse generalisation check only, because full scenes are intrinsically harder than chunks, so the ratio mixes difficulty with overfitting and we read it cautiously. We assess semantic structure qualitatively by principal-component visualisations of the per-primitive feature space and quantitatively by clustering and linear-probe measures on those features. Stage~2 generation and completion are evaluated by generation-quality and fidelity metrics for which we likewise leave space.

        \chapter{Results}
\label{sec:results}

We organize results by research question. \autoref{sec:res_scaling} addresses RQ1 (scaling and reconstruction), \autoref{sec:res_identifiability} addresses RQ2 (identifiability, ordering, and sampling), \autoref{sec:res_semantic} addresses RQ3 (semantic structure), and \autoref{sec:res_stage2} addresses RQ4 (generation and completion). Where the large-scale model is still training at the time of writing, we state the experimental design and the qualitative findings that are already established, and leave clearly marked space for the final quantitative numbers. Tables whose cells read ``--'' are placeholders to be filled from the final logs; their structure and captions fix what each result will report.

\paragraph{A note on the reported numbers.} The small-scale disentanglement and semantic-structure numbers in \autoref{sec:res_semantic} are taken from our exploratory experiments and are included to show the structure of the analysis; they should be re-confirmed against the final experiment logs before submission. The large-scale reconstruction, fidelity, and Stage 2 numbers are left as placeholders.

% =====================================================================
\section{Scaling and Reconstruction (RQ1)}
\label{sec:res_scaling}
% =====================================================================

\paragraph{Small scale converges; naive scaling does not.} With the inherited encoder and decoder, a small training set on the order of a few hundred scenes converges well: reconstruction error falls smoothly and the model generalizes to held-out scenes, reproducing the regime in which Can3Tok was validated \citep{gao2025can3tok}. When we scaled the training set to several thousand heterogeneous scenes, the same configuration failed to converge to a comparable reconstruction, which is the central empirical problem of this thesis and the motivation for the encoder, decoder, objective, and data-pipeline changes of \autoref{sec:methodology}. \autoref{tab:scaling} reports the reconstruction of the final large-scale model against the small-scale reference and an inherited-architecture baseline that we are training for comparison.

\begin{table}[h]
  \centering
  \caption[Stage 1 reconstruction at scale]{\textbf{Stage 1 reconstruction.} Held-out reconstruction error and render fidelity for the small-scale reference, the inherited-architecture baseline trained at scale, and our final model. Lower is better for error and LPIPS; higher is better for PSNR and SSIM. \emph{[TODO: fill from final logs.]}}
  \label{tab:scaling}
  \begin{tabular}{lccccc}
    \textbf{Model} & \textbf{Pos} $\downarrow$ & \textbf{Cov} $\downarrow$ & \textbf{PSNR} $\uparrow$ & \textbf{SSIM} $\uparrow$ & \textbf{LPIPS} $\downarrow$ \\
    \hline
    Small scale (ref.)            & -- & -- & -- & -- & -- \\
    Inherited arch., at scale     & -- & -- & -- & -- & -- \\
    Ours (structured + local)     & -- & -- & -- & -- & -- \\
  \end{tabular}
\end{table}

\paragraph{Reconstruction quality of the final model.} \autoref{fig:recon} shows qualitative reconstructions of held-out scenes by the final model, alongside the ground-truth renders, to be inserted once training completes.

\begin{figure}[h]
  \centering
  \includegraphics[width=0.92\linewidth]{example-image}
  \caption[Qualitative reconstruction]{\textbf{Qualitative reconstruction.} Ground-truth (top) and reconstructed (bottom) renders of held-out scenes from the final model. \emph{[FIGURE TODO: insert paired ground-truth/reconstruction renders exported from the reconstruction PLYs.]}}
  \label{fig:recon}
\end{figure}

% =====================================================================
\section{Identifiability, Ordering, and Sampling (RQ2)}
\label{sec:res_identifiability}
% =====================================================================

This section isolates the factors that we argue are responsible for the scaling behaviour: the serialization order, the number of primitives sampled, and the locality of the encoder.

\paragraph{Serialization order.} Replacing the opacity rank with a space-filling-curve order makes the element-wise reconstruction target spatially stable and is, in our experience, the single change that most improved large-scale convergence, consistent with the locality argument for the Hilbert curve over the Z-order curve \citep{wu2024ptv3}. We note, however, that ordering alone improved convergence but did not by itself close the gap to small-scale quality, which is why it is necessary but not sufficient and is combined with the sampling and locality changes below. \autoref{tab:ordering} reports the ablation over no ordering, Z-order, and Hilbert order.

\begin{table}[h]
  \centering
  \caption[Effect of serialization order]{\textbf{Serialization order.} Reconstruction error under opacity rank (no spatial order), Z-order (Morton), and Hilbert order, all else equal. \emph{[TODO: fill from logs. Expected trend: Hilbert $<$ Morton $<$ opacity in error.]}}
  \label{tab:ordering}
  \begin{tabular}{lcc}
    \textbf{Order} & \textbf{Recon error} $\downarrow$ & \textbf{Gen. gap} $\rightarrow 1$ \\
    \hline
    Opacity rank (none) & -- & -- \\
    Z-order (Morton)    & -- & -- \\
    Hilbert             & -- & -- \\
  \end{tabular}
\end{table}

\paragraph{Number of sampled primitives.} We find that, with pre-fitted and frozen splats, sampling a smaller budget of primitives (on the order of $10$k) converges markedly better than a larger budget (on the order of $40$k), the opposite of what the larger budget achieves in the original setting where the splats are jointly optimized during sampling \citep{gao2025can3tok}. We read this as a consequence of the within-block placement gauge: a larger budget packs more primitives per latent token, increasing the nuisance that the token must encode. \autoref{tab:sampling} reports the ablation.

\begin{table}[h]
  \centering
  \caption[Effect of the number of sampled primitives]{\textbf{Sampling budget.} Reconstruction error as a function of the number of primitives sampled per scene, with frozen splats. \emph{[TODO: fill from the $10$k and $40$k training runs.]}}
  \label{tab:sampling}
  \begin{tabular}{lcc}
    \textbf{Primitives per scene} & \textbf{Recon error} $\downarrow$ & \textbf{Converged?} \\
    \hline
    $\sim 10$k & -- & -- \\
    $\sim 40$k & -- & -- \\
  \end{tabular}
\end{table}

\paragraph{Locality of the encoder and the generalization gap.} The change that, in our experiments, most directly closed the train-to-test generalization gap at scale was replacing the globally-mixed bottleneck with the spatially local per-token encoder of \autoref{sec:method_local}. With the global encoder the held-out error substantially exceeded the training error (a gap on the order of five to eight times); with the structured local encoder the gap closed to approximately unity, indicating that the latent had become compositional rather than a memorized global code. \autoref{tab:encoder} reports the ablation between the global and local encoders, including the window-scope control.

\begin{table}[h]
  \centering
  \caption[Global versus local encoder]{\textbf{Encoder locality and generalization.} Training error, held-out error, and their ratio for the global encoder and the spatially local encoder. A ratio near one indicates no overfitting. \emph{[TODO: fill from logs. Established qualitative finding: gap $\approx 5$--$8\times$ (global) $\rightarrow \approx 1\times$ (local).]}}
  \label{tab:encoder}
  \begin{tabular}{lccc}
    \textbf{Encoder} & \textbf{Train} $\downarrow$ & \textbf{Held-out} $\downarrow$ & \textbf{Gap} $\rightarrow 1$ \\
    \hline
    Global bottleneck      & -- & -- & $\sim 5$--$8\times$ \\
    Local (window $\omega$) & -- & -- & $\sim 1\times$ \\
  \end{tabular}
\end{table}

\paragraph{Objective.} The gauge-invariant covariance loss and the permutation-invariant set loss are evaluated by the gauge-invariant covariance error and by the colour and orientation fidelity of reconstructions, since the raw quaternion error is gauge-blind and cannot register their effect. \autoref{tab:objective} reports the ablation.

\begin{table}[h]
  \centering
  \caption[Effect of the reconstruction objective]{\textbf{Reconstruction objective.} Gauge-invariant covariance error and colour fidelity under the element-wise objective, the gauge-invariant covariance objective, and the permutation-invariant set objective. \emph{[TODO: fill from logs.]}}
  \label{tab:objective}
  \begin{tabular}{lcc}
    \textbf{Objective} & \textbf{Cov. error} $\downarrow$ & \textbf{Colour fidelity} $\uparrow$ \\
    \hline
    Element-wise ($L_2$)        & -- & -- \\
    + gauge-invariant covariance & -- & -- \\
    + permutation-invariant set  & -- & -- \\
  \end{tabular}
\end{table}

% =====================================================================
\section{Semantic Structure (RQ3)}
\label{sec:res_semantic}
% =====================================================================

\paragraph{Per-primitive semantic clustering.} The per-primitive contrastive objective organizes the per-primitive feature space so that primitives of the same category cluster together. \autoref{fig:pca} shows principal-component projections of the per-primitive features, in which object regions emerge as coherently coloured clusters, which is the qualitative evidence that the latent carries semantic structure rather than only geometry. These visualizations were already obtained in the small-scale experiments and will be regenerated for the final model.

\begin{figure*}[t]
  \centering
  \includegraphics[width=0.95\linewidth, height=0.24\textheight]{example-image}
  \caption[Semantic feature visualization]{\textbf{Per-primitive semantic structure.} Principal-component projection of the per-primitive contrastive features, coloured by the leading components, for several scenes. Coherent regions correspond to object categories. \emph{[FIGURE TODO: insert the PCA feature PLY renders. Source: the semantic-feature and z\_s visualizations.]}}
  \label{fig:pca}
\end{figure*}

\paragraph{Colour via the global shape embedding.} Routing the per-scene mean colour through a head on the global shape embedding, and decoding colour as a residual about it, improved colour fidelity in the small-scale experiments by separating the coarse colour statistic (carried globally) from the fine per-primitive variation (decoded locally). This is reflected in the colour-fidelity column of \autoref{tab:objective} and in the qualitative reconstructions.

\paragraph{Disentanglement and latent structure.} In the small-scale regime we measured the effect of the disentangled latent and its objectives on (i) category clustering in the layout subspace, (ii) the degree to which swapping the layout subspace across scenes leaves geometry intact (a measure of disentanglement), and (iii) the stability of the layout code under partial observation (a proxy for completion readiness). \autoref{tab:disentangle} reproduces these exploratory results; they indicate that the disentanglement objectives produced a real separation (swapping the layout subspace disrupted geometry far less than a cross-scene baseline) and that a deterministic layout code was the most stable under partial observation. These numbers are from the exploratory phase and are to be confirmed.

\begin{table*}[t]
  \centering
  \caption[Small-scale disentanglement and conditioning ablation]{\textbf{Small-scale latent-structure ablation (exploratory, to be confirmed).} Layout-subspace category clustering (higher is better), geometry disruption under a cross-scene layout swap as a fraction of the cross-scene baseline (lower indicates better disentanglement), layout-code stability under a partial observation (higher is better), and fraction of prior samples decoding to valid geometry. Reproduced from the exploratory experiment notes. \emph{[TODO: confirm against final logs; extend to the large-scale model.]}}
  \label{tab:disentangle}
  \begin{tabular}{lcccc}
    \textbf{Variant} & \textbf{Layout clustering} $\uparrow$ & \textbf{Swap disruption} $\downarrow$ & \textbf{Partial-obs.\ stability} $\uparrow$ & \textbf{Valid \%} $\uparrow$ \\
    \hline
    Baseline (undifferentiated)        & 1.05 & 6.5\%  & --    & 100\% \\
    Disentangled                        & 0.99 & 0.27\% & 0.26  & 94\%  \\
    Disentangled + structured split     & 0.97 & 0.23\% & 0.20  & 81\%  \\
    + per-primitive contrastive          & 0.98 & 0.08\% & 0.37  & 95\%  \\
    Cross-attention conditioning         & 1.24 & 0.34\% & 0.76  & 97\%  \\
  \end{tabular}
\end{table*}

\paragraph{Semantic feature quality.} \autoref{tab:semquality} reports clustering and linear-probe measures of the per-primitive feature space for the contrastive variants, again from the exploratory phase. They show modest but real separability of categories in the learned features, which the principal-component visualizations corroborate.

\begin{table}[h]
  \centering
  \caption[Semantic feature quality]{\textbf{Per-primitive feature quality (exploratory, to be confirmed).} Between/within-class separation and linear-probe accuracy of the per-primitive features. \emph{[TODO: confirm; recompute for the final model.]}}
  \label{tab:semquality}
  \begin{tabular}{lcc}
    \textbf{Measure} & \textbf{Contrastive} & \textbf{Contrastive + pooled} \\
    \hline
    Between/within separation $\uparrow$ & 0.44 & 0.41 \\
    Linear-probe accuracy $\uparrow$     & 26.0\% & 26.1\% \\
  \end{tabular}
\end{table}

% =====================================================================
\section{Generation and Completion (RQ4)}
\label{sec:res_stage2}
% =====================================================================

\paragraph{Generation.} \autoref{fig:generation} will show unconditional and conditional scene generations from the Stage 2 model decoded by the frozen Stage 1 decoder, with generation-quality and fidelity metrics in \autoref{tab:stage2}.

\begin{figure}[h]
  \centering
  \includegraphics[width=0.92\linewidth]{example-image}
  \caption[Scene generation]{\textbf{Stage 2 scene generation.} Scenes sampled from the latent flow-matching model and decoded by the frozen Stage 1 decoder. \emph{[FIGURE TODO: insert generated-scene renders once Stage 2 is trained.]}}
  \label{fig:generation}
\end{figure}

\begin{table}[h]
  \centering
  \caption[Stage 2 generation and completion metrics]{\textbf{Stage 2 metrics.} Generation quality and, for completion, reconstruction fidelity against held-out full scenes from partial observations. \emph{[TODO: fill once Stage 2 is trained. Completion is left as future work; see \autoref{sec:discussion}.]}}
  \label{tab:stage2}
  \begin{tabular}{lccc}
    \textbf{Task} & \textbf{Quality} $\uparrow$ & \textbf{Fidelity} $\uparrow$ & \textbf{Coverage} \\
    \hline
    Unconditional generation & -- & -- & -- \\
    Conditional generation   & -- & -- & -- \\
    Scene completion         & -- & -- & -- \\
  \end{tabular}
\end{table}

\paragraph{Completion.} As stated in \autoref{sec:method_stage2} and discussed in \autoref{sec:discussion}, scene completion is specified and set up but not brought to a finished, evaluated result, because the effort available was absorbed by the Stage 1 scaling problem. \autoref{fig:completion} is reserved for completion results.

\begin{figure}[h]
  \centering
  \includegraphics[width=0.92\linewidth]{example-image}
  \caption[Scene completion]{\textbf{Scene completion (reserved).} Partial input scan, completion, and ground truth. \emph{[FIGURE TODO / FUTURE WORK: insert if completion is brought to a result.]}}
  \label{fig:completion}
\end{figure}
        \chapter{Discussion}
\label{chap:discussion}

This chapter steps back from the numbers and asks what the seven-variant ablation and the Stage~2 results actually support. I revisit each research question in turn, then discuss what the frozen-splat regime is doing to the loss landscape, then the position compression I observed in Stage~2 and what I tried to fix it, and finally the limitations I am aware of and that constrain the conclusions one can draw. I write this chapter in the spirit of an honest project post-mortem rather than a victory lap: the architectural changes I propose are large quantitative improvements, and the Stage~2 generator does work, but the system as a whole is not yet at a quality I would call photorealistic, and I want to be precise about where the residual difficulty lies.

%% ──────────────────────────────────────────────────────────────────────────
\section{Revisiting the Research Questions}
\label{sec:disc_rqs}

\paragraph{RQ1 (Ordering).} The result is positive but with one qualifier. Hilbert serialisation in a fixed canonical frame substantially improves validation reconstruction (V1 $\to$ V2: $99 \to 77$ on the full-scene validation $\ell_2$, Table~\ref{tab:7_summary}). The improvement is entirely a generalisation effect; on the training set, the V1 architecture can also bring the loss down, but it does so by storing per-scene context in the bottleneck. Imposing a canonical order on the target makes the position component of the loss into a learnable function, so the network can use the same parameters across scenes instead of memorising a per-scene mapping. The qualifier is that the ordering only buys what it buys; pairing it with a position scaffold gets a much larger drop (V2 $\to$ V3: $77 \to 45$), so the order is the necessary precondition rather than the full solution.

\paragraph{RQ2 (Architecture and generalisation).} The result is also positive but with a clearer mechanism than I expected. Replacing the flat $777$M-parameter decoder with a $1.6$M-parameter token-local MLP \emph{by itself} does not close the gap (V1 $\to$ V4: stays at $99$). Adding a local encoder \emph{by itself} would not be testable cleanly because the local encoder gives a per-token latent that the flat decoder cannot use coherently. The two changes only work together. When both are on (V5), reconstruction improves to $79$ with a $500\times$ reduction in decoder parameters; when the Hilbert order plus scaffold are added on top (V6), reconstruction drops to about $47$, comparable to V3's $45$ but with two orders of magnitude fewer decoder parameters. The mechanism is that the local encoder and the token-local decoder are two halves of one structural decision: latent token $k$ is a code for spatial region $k$, and the decoder MLP only has to model what is in that region. This compositional structure is what generalises. The qualitative figure (Figure~\ref{fig:recon}) shows that the visual jump happens between V3 and V6 even though the aggregate metric only moves a few points, which is consistent with this reading.

\paragraph{RQ3 (Semantic structure).} The 300-scene experiments (\S\ref{sec:results_300}) show that per-Gaussian InfoNCE based on ScanNet-72 labels can be added at no measurable cost to reconstruction, while producing per-Gaussian features that cluster by category (Figure~\ref{fig:pca}). The layout/geometry split (Strategy~B1's deterministic-layout idea, lifted into the VAE form for the large-scale runs) gives a more stable layout code under partial-scan masking. At scale (V6 vs.\ V3) the disentangled and per-Gaussian-supervised variant matches the un-disentangled variant on reconstruction, and provides the schema (a 16-token $z_s$ plus a 512-token $z_g$) that the hierarchical Stage~2 generator needs. The Stage~2 results confirm that this schema is usable: \texttt{LayoutDiT} reaches a low validation loss on the 16-token layout latent and \texttt{GeometryDiT512} conditioned on it reaches a comparably low loss on the 512-token geometry latent. The qualifier here is that semantic structure was not the limiting factor for reconstruction at any scale studied, so the disentanglement is best read as a Stage-2-enabling design choice rather than a Stage-1 reconstruction win.

\paragraph{Stage 2 generation versus completion.} The qualitative Stage~2 result on V6 (Figure~\ref{fig:stage2_exp6}) is honestly mixed. Unconditional generation reproduces the structural pattern of the Stage~1 latent well; the shape of the generated point cloud closely matches the Stage~1 reconstruction of a held-out scene. Completion under $40\%$ observation recovers the coarse layout but is more diffuse than the reference, with the high-density region spreading along the diagonal of the scene. Two interpretations are consistent with this: completion is a constrained inverse problem on an imperfect latent, and the Stage~1 decoder is particularly sensitive to residual error in the position channel, so any inaccuracy in the completed latent shows up as positional smearing in the rendered output. I view the result as evidence that the Stage~2 architecture and training pipeline are sound (the flow matching converges, generation works), and that the residual limitation lies in the precision of the inferred latent for unobserved regions.

%% ──────────────────────────────────────────────────────────────────────────
\section{What the Frozen-Splat Regime Does to the Loss Landscape}
\label{sec:disc_frozen}

The contrast between the Can3Tok original demonstration (a few hundred scenes, splat-fitting jointly optimised with the encoder) and the SceneSplat-7K setting (thousands of scenes, splats pre-fitted and frozen) deserves its own paragraph because it is, in my reading, the largest residual obstacle to closing the photorealism gap. In Can3Tok's joint regime, the splat-fitting process can subtly co-adapt to whatever per-token tokenisation the encoder ends up using. The decoder slot $i$ has implicit influence on what gets put into the input position $i$. In the pre-fitted regime, this co-adaptation is broken: the splats are fixed by a renderer-quality objective that has no knowledge of the autoencoder, so the network must absorb whatever within-region arrangement the renderer produced.

This explains why the same architecture that worked at $300$ scenes in the joint setting fails to generalise at $6{,}000$ in the pre-fitted setting. The 300-scene exploration in this thesis (also pre-fitted) does not contradict this; at 300 scenes the network is still in the memorisation regime where per-scene context is cheap to store, and the failure mode of the pre-fitted setting is not yet active. The seven-variant ablation is in some sense a study of how to make a slot-aligned loss work when the data-side cooperation that the joint setting silently relied on is removed.

It also explains the otherwise surprising finding that the position scaffold (V3) gives a much larger improvement than ordering alone (V2). With ordering alone, the network still has to predict the absolute coordinate at slot $i$, which is hard for the same reason: within a Hilbert block the absolute position varies room by room. With the scaffold, the network only has to predict a bounded local offset relative to a token-specific anchor that is itself constrained to be the block centroid. The hard work is delegated to the anchor head, which has a structurally easy job, and the offset head solves a simple local-deformation problem. The two-stage decomposition matches what the pre-fitted regime can actually identify from data.

%% ──────────────────────────────────────────────────────────────────────────
\section{Position Compression and What I Tried}
\label{sec:disc_compression}

Across all Stage~2 diagnostics, the positional output of the Stage~1 decoder, when fed Stage~2 latents, lies in a smaller range than the canonical $\pm 10$\,m training frame. V6 generates positions in roughly $\pm 2$\,m; V7 generates them in approximately $\pm 0.7$\,m. Real validation scenes from the same set occupy roughly $\pm 5$ to $\pm 9$\,m in the canonical frame. The compression is consistent across epochs, across the layout, geometry, and completion stages, and across the four scenes saved at each evaluation checkpoint.

I diagnosed the source carefully. The flow-matching cosine similarity between predicted and target velocity sits above $0.96$ for V6 geometry from epoch 50 onwards, which means the DiT is learning the velocity field correctly: the latent it produces matches the latent the encoder produces. The compression therefore is not introduced by Stage~2. Encoding a real validation scene and decoding it directly through the frozen Stage~1 decoder reproduces the same compressed range, which places the issue in the Stage~1 decoder output rather than in any Stage~2 module. The anchor-relative position decoding (Eq.~\eqref{eq:anchor_relative}) bounds the local offset to $\rho \tanh(o_i)$ with $\rho$ a learnable global scale; my best reading is that the network has learned a small $\rho$ during Stage~1 training because the slot-aligned $\ell_2$ loss rewards conservative predictions when the anchor is approximately correct, and the bias persists through the frozen decoder used at Stage~2.

What I tried to fix it during the project:
\begin{itemize}
\item \textbf{Loosening the anchor offset bound.} Re-initialising $\rho$ at a larger value (4.0 instead of 2.0) for a partial re-training run did not change the qualitative compression. The network re-converges to a similar effective scale.
\item \textbf{Diagnostic on the decoded coordinates.} I verified that the Stage~1 decoder receives a latent in the expected range during Stage~2 sampling and that the inverse normalisation in the PLY writer is applied symmetrically to GT and generated PLYs, ruling out a writer-side bug.
\item \textbf{Comparing V6 and V7.} The compression is worse in V7 (no scaffold) than V6 (scaffold), which suggests the scaffold is doing the right thing structurally but the anchor centroids themselves are not at the right absolute scale.
\end{itemize}

What I would do next if the project continued: re-train Stage~1 with an explicit positional-range loss that penalises the output for being narrower than the input, or, more cleanly, parameterise $\rho$ per Hilbert block rather than as a global scalar so that high-extent regions of the scene can use a wider offset range. Both are tractable but were outside the time budget. The honest framing is that Stage~2 demonstrates the latent generative pipeline works in principle, while Stage~1 still has a positional-range bug that holds back the final rendered quality.

%% ──────────────────────────────────────────────────────────────────────────
\section{Limitations}
\label{sec:disc_limitations}

I list the limitations I am aware of, ordered roughly from most consequential.

\paragraph{Stage 1 has not converged to a photorealistic level.} Even the best Stage~1 variant (V6, validation $\ell_2 = 47.5$) produces reconstructions that are recognisable as ``a room with multiple distinct objects'' but are far from a per-scene 3DGS fit in visual quality (Figure~\ref{fig:recon}). The aggregate metric and the qualitative figure both point in the same direction. The thesis demonstrates that the architectural changes I propose move the system from a complete-failure regime (V1) to a regime where the scene structure becomes visible (V6), but the system is not yet at a quality that competes with per-scene 3DGS reconstructions or with the original Can3Tok demonstrations at the 300-scene scale.

\paragraph{Position compression in the decoder output.} As described in \S\ref{sec:disc_compression}, the Stage~1 decoder produces positions in a smaller range than the canonical training frame, and this propagates to all Stage~2 outputs. I have characterised the source but the fix was not in scope for this thesis.

\paragraph{No rendering-quality metrics.} Every number in Chapter~\ref{chap:results} is on the same per-primitive $\ell_2$ residual that is used during training (Stage~1) or the rectified-flow velocity MSE (Stage~2). I do not report PSNR, SSIM, or LPIPS on rendered images of the reconstructed scenes. The per-primitive $\ell_2$ is a useful relative metric for comparing variants but is not directly tied to perceptual quality. Adding rendered-image evaluation is the next concrete step.

\paragraph{Stage 2 completion is partial.} Generation works on the V6 latent (Figure~\ref{fig:stage2_exp6}b), reproducing the structural pattern of the Stage~1 latent. Completion (panel d) is recognisably scene-shaped but more diffuse than the reference. The training loss converges and the flow diagnostics are clean, so I attribute the qualitative gap to the constrained-inversion nature of completion combined with the Stage~1 decoder's sensitivity to small latent errors in the position channel.

\paragraph{Fixed-$N$ assumption.} I work with a fixed $40{,}000$ Gaussians per scene, sampled by top-opacity ranking. Real 3DGS scenes have variable counts (often in the hundreds of thousands), and the top-opacity ranking biases the input distribution towards opaque furniture surfaces. A density-uniform farthest-point sampling would give a more spatially balanced input and is a clear next experiment.

\paragraph{Indoor-only.} SceneSplat-7K is indoor only. The Hilbert canonical frame (radius $10$\,m) is sized for indoor extents. Outdoor scenes with unbounded depth would need either a different normalisation (log-depth contracted) or a tiled-scene approach. I have not tested either.

\paragraph{No comparison to a Can3Tok-style joint training at scale.} A fair test of the frozen-vs-joint hypothesis (\S\ref{sec:disc_frozen}) would be to also train a joint-splat variant at the same scale and compare. Doing this is computationally expensive (splat fitting per scene is not free) and was out of scope for this thesis.

\paragraph{Hyperparameter sweep is shallow.} For each of the seven variants I have a single training run on a single seed. Variance estimates would be informative, but each run takes a couple of days on $2 \times $H100 and the budget did not allow for repeats. The trends between variants are large enough relative to plausible seed variance that I am confident in the qualitative ordering, but precise effect sizes should be treated as point estimates.

\paragraph{ScanNet-72 labels do not transfer cleanly across sub-datasets.} I disable per-Gaussian semantic supervision for the ARKitScenes and ScanNet++ portions of the training corpus because their taxonomies do not match ScanNet-72 exactly. This is a pragmatic choice; a unified label space (or a multi-taxonomy head) would let the per-Gaussian supervision act on every training scene.

\paragraph{Disentanglement is structural, not identifiable.} The layout/geometry split in V6/V7 relies on auxiliary heads (scene-mean colour, scene-label distribution, per-category centroids), cross-reconstruction, and the orthogonality regulariser. These are structural inductive biases. As Locatello et al.\ \citep{locatello2019challenging} point out, unsupervised disentanglement is not identifiable; my split is identifiable only insofar as the auxiliary supervision pins down what $z_s$ is supposed to contain. The supervision could be stronger or weaker without my having a principled way to choose.

None of these limitations invalidates the central finding of the thesis (the joint role of Hilbert ordering, anchor-relative decoding, and the local encoder plus token-local decoder pair in scaling the Stage~1 VAE, and the schema-aware Stage~2 generator that follows from it). They do delimit the scope within which that finding is presently supported.
        \chapter{Conclusion}
\label{chap:conclusion}

This thesis asked what it takes to scale a scene-level 3D Gaussian Splatting variational autoencoder from a few hundred training scenes (the regime of Can3Tok \citep{gao2025can3tok}) to several thousand pre-fitted heterogeneous indoor scenes (the regime made available by SceneSplat-7K \citep{li2025scenesplat}), and what the design implications are for a downstream flow-matching generator. I conclude by returning to the three research questions and listing the next concrete steps.

%% ──────────────────────────────────────────────────────────────────────────
\section{Answers to the Research Questions}
\label{sec:conc_answers}

\paragraph{RQ1.} Yes, imposing a canonical spatial order on the Gaussian set restores convergence at scale, but only when paired with a decoder that can use it. A fixed-frame Hilbert serialisation (\S\ref{sec:method_hilbert}) drops validation $\ell_2$ from 99 to 77 with no other change (V1 $\to$ V2). Adding anchor-relative position decoding on top of the ordering drops it further to 45 (V2 $\to$ V3). The ordering and the anchor decomposition together turn an unlearnable slot-aligned objective into a learnable one.

\paragraph{RQ2.} Yes, a windowed local encoder plus a shared per-token decoder closes most of the train-to-test gap. The two components are interlocking: the local encoder needs the token-local decoder to consume per-region features cleanly, and the token-local decoder needs the local encoder to feed it the right kind of context. Together, the decoder parameter count drops from about 777M to about 1.6M (a $500\times$ reduction) while validation reconstruction reaches a level comparable to the global-decoder run.

\paragraph{RQ3.} Yes, per-Gaussian semantic supervision and a layout/geometry latent split can be added without degrading reconstruction. The 300-scene exploratory study (Strategy~B1 with Pool+pgNCE-style InfoNCE, \S\ref{sec:results_300}) confirms that per-Gaussian features cluster by ScanNet-72 category and that a deterministic layout code is stable under partial-scan masking. At scale, the V6/V7 variants match V3 on reconstruction while producing the split-schema latent that the hierarchical Stage~2 generator (\S\ref{sec:method_stage2}) needs.

\paragraph{Main question.} The combination of fixed-frame Hilbert serialisation, anchor-relative position decoding, and a local-encoder plus token-local-decoder pair is what scales the Stage~1 VAE from hundreds to several thousand pre-fitted indoor scenes. Each ingredient on its own is insufficient; the three together turn an architecture that fails to converge at scale into one that reconstructs held-out scenes at a useful level with a 500$\times$ smaller decoder. Layered on top, a layout/geometry split with per-Gaussian semantic supervision is a viable schema for the downstream generator.

%% ──────────────────────────────────────────────────────────────────────────
\section{Future Work}
\label{sec:conc_future}

The next milestones for the project follow directly from the limitations in \S\ref{sec:disc_limitations}.

\begin{enumerate}
\item \textbf{Finish Stage~2 end-to-end.} Generate samples from the trained Stage~2 models for variants 5, 6, and 7, render them, and report PSNR, SSIM, and LPIPS against held-out ground truth. The same checkpoints support the completion task: mask 30--80\% of the input Hilbert blocks of a held-out scene and recover the masked regions through Stage~2.
\item \textbf{Density-uniform input sampling.} Replace the top-opacity ranking of the 40{,}000 input Gaussians with farthest-point sampling on positions, weighted by opacity. This should give a more spatially uniform input distribution and reduce the bias towards opaque furniture surfaces.
\item \textbf{Rendering-quality evaluation throughout.} Add rendered-image metrics (PSNR/SSIM/LPIPS) on the reconstructed scenes for all seven variants, not just the generative outputs. This gives the perceptual companion metric to the per-primitive $\ell_2$.
\item \textbf{Unified label space.} Bring ARKitScenes and ScanNet++ under a unified per-Gaussian taxonomy (a multi-taxonomy classification head, with per-dataset masks) so the InfoNCE supervision can be active on the full training corpus.
\item \textbf{Outdoor scenes.} Test the same architecture on outdoor 3DGS data with a contracted-depth normalisation. The Hilbert serialisation and the anchor scaffold should both generalise, but the canonical extent has to change.
\item \textbf{Joint-splat ablation.} Re-train one of the Stage~1 variants jointly with the Gaussian fitting at scale and compare to the same variant on the pre-fitted SceneSplat-7K splats. This is the cleanest test of the frozen-splat hypothesis (\S\ref{sec:disc_frozen}).
\item \textbf{Multiple seeds.} Re-run V3, V5, V6, V7 on at least three seeds each to get variance estimates on the validation $\ell_2$ comparisons.
\end{enumerate}

The codebase, the SLURM job scripts, the YAML configs for each of the seven variants, and the Stage~2 model code are available at \repo. The Stage~1 checkpoints used to produce Table~\ref{tab:7_summary} are released alongside, so the next person can pick up directly from the Stage~2 milestone without retraining.

\bibliographystyle{ACM-Reference-Format}
\bibliography{bibliographies/references}

% \newpage
\appendix
        % You can choose whether you prefer a single or double column appendix.
% \onecolumn

\chapter{Implementation Details}
\label{sec:apx:first_appendix}

This appendix records the architecture, the training configuration, the data, the gauge-invariant objective, and the evaluation metrics in enough detail to reproduce the experiments. As stated in \autoref{sec:methodology}, the thesis is intended to be readable without it.

% ---------------------------------------------------------------------
\section{Model Architecture and Dimensions}
% ---------------------------------------------------------------------

The encoder is a cross-attention Perceiver. A set of learned query tokens cross-attends to the input Gaussian set and is followed by self-attention; one query summarizes the scene globally and the remaining queries carry per-region geometry. The transformer width is $384$, with $12$ attention heads, $6$ encoder self-attention layers, and $12$ decoder layers, and the Fourier position embedding uses $8$ frequency bands. The latent has $512$ tokens of dimension $32$, so the latent carries $512 \times 32$ values; the per-token dimension is the capacity knob, with the token count pinned to the number of Hilbert blocks.

In the structured, spatially local configuration used for the large-scale model, the encoder is windowed: each geometry query attends only to a local window of Hilbert blocks around its own block (a half-width of one block in the reported run), with a single global query attending to all primitives, so that latent token $k$ encodes the local geometry of block $k$. The decoder is token-local: a shared per-token network decodes each token's slice of the primitives, replacing the flat decoder that otherwise dominates the parameter count. Position is decoded anchor-relative in the spirit of Scaffold-GS \citep{lu2024scaffold}: each primitive's position is an anchor (the per-block centroid) plus a bounded local offset, with the bound a single learnable global scale initialized at $2.0$. A Fourier decoder position embedding over a fixed scaffold grid is added to the decoder tokens. The latent-disentanglement and cross-attention conditioning variants of \autoref{sec:method_latent} share this encoder and differ only in how the latent is split and how the decoder is conditioned.

% ---------------------------------------------------------------------
\section{Training Configuration}
% ---------------------------------------------------------------------

The large-scale model is trained on two H100 GPUs in \texttt{bfloat16} mixed precision with distributed data parallelism. The batch size is between $64$ and $90$ depending on the configuration. Optimization uses AdamW with a learning rate of $10^{-4}$, weight decay $10^{-2}$, a linear warmup of $300$ steps, and a cosine schedule with warm restarts. The KL weight is fixed at $10^{-5}$, deliberately small so that the latent is lightly regularized and the second-stage generator can model its distribution. The reconstruction objective in the reported run combines the permutation-invariant set loss (entropic optimal transport with a regularization of $0.05$ and $50$ iterations, with unit position and opacity weights) with the gauge-invariant covariance term (the co-diagonalization Bures surrogate, with a small isotropic regularizer, fifteen Newton-Schulz iterations, and a shape weight of $5.0$). Colour is decoded as a residual about a per-scene mean predicted from the global embedding, with a mean-colour weight of $0.5$ and a colour weight of $2.0$. The per-primitive contrastive objective uses the hidden-feature variant with a weight of $0.4$, a temperature of $0.10$, and a balanced subsample of primitives per scene. The Gaussian set is serialized along a fixed-frame Hilbert curve over the canonical radius, and primitives are selected by opacity rank. The sampling-budget ablation of \autoref{tab:sampling} varies the number of sampled primitives between roughly $10$k and $40$k with all else fixed.

Each Gaussian is represented by fourteen values: a three-dimensional position, a three-dimensional colour (a residual when the colour-residual decomposition is used), a scalar opacity, a three-dimensional scale, and a four-dimensional quaternion in $(w, x, y, z)$ order. Scale is decoded in linear space, which the covariance objective requires in order to assemble $\Sigma = R\,\mathrm{diag}(s^2)\,R^\top$.

% ---------------------------------------------------------------------
\section{Configurations of the Ablation Runs}
% ---------------------------------------------------------------------

\autoref{tab:apx_runs} lists the configurations behind the reconstruction ladder of \autoref{tab:ladder}. All runs share the same training mixture, the $512$-token latent of width $32$, the Hilbert serialisation, opacity sampling at $40{,}000$ primitives per scene, the colour-residual decomposition, and the Fourier decoder position embedding; they differ only in the encoder, the decoder, the anchor-relative position decoding, and the latent split. ``Reorder'' marks the runs that additionally apply the Morton-index reordering pass on top of the Hilbert blocks. The position error is the best held-out full-scene value from \autoref{tab:ladder}; the entry for the reorder-only run is left blank because its summary was not available when this table was assembled and should be filled from its log.

\begin{table*}[t]
  \centering
  \caption[Configurations of the ablation runs]{\textbf{Ablation-run configurations.} The seven Stage~1 runs behind \autoref{tab:ladder}, by encoder (global bottleneck or windowed local), decoder (flat or token-local), anchor-relative position decoding, latent organisation (undifferentiated, structured per-token without a semantic split, or disentangled into $z_s$ and $z_g$), and the optional Morton reorder. Position error is the best held-out full-scene value.}
  \label{tab:apx_runs}
  \begin{tabular}{llllll c}
    \textbf{Run} & \textbf{Encoder} & \textbf{Decoder} & \textbf{Anchor} & \textbf{Latent} & \textbf{Reorder} & \textbf{Pos} $\downarrow$ \\
    \hline
    E1 & global & flat        & no  & undifferentiated      & no  & $5386$ \\
    E2 & global & flat        & no  & undifferentiated      & yes & -- \\
    E3 & global & flat        & yes & undifferentiated      & yes & $543$ \\
    E4 & global & token-local & no  & undifferentiated      & no  & $5372$ \\
    E5 & local  & token-local & no  & structured per-token  & no  & $2415$ \\
    E6 & local  & token-local & yes & disentangled ($z_s,z_g$) & yes & $\mathbf{451}$ \\
    E7 & local  & token-local & no  & disentangled ($z_s,z_g$) & yes & $541$ \\
  \end{tabular}
\end{table*}

% ---------------------------------------------------------------------
\section{Data and Preprocessing}
% ---------------------------------------------------------------------

Training uses SceneSplat-7K \citep{li2025scenesplat}. The main set is a collection of indoor chunks expressed in a global per-scene coordinate frame, obtained by combining all chunks of a scene to compute one normalization factor so that all chunks of a room share a coordinate system. Additional full scenes from several indoor sources are added, each normalized per scene to the same canonical radius; per-primitive semantic labels are present only on the chunk subset, so the semantic objective is active there and disabled on the additional sources, which avoids a cross-source label-taxonomy collision. Validation uses held-out full scenes for the primary metric and a disjoint set of held-out chunks for the distribution-gap diagnostic; the chunk split is constructed so that the training and held-out chunks never overlap.

Positions are normalized into a canonical sphere of fixed radius. Colours are scaled to the unit interval and, when the residual decomposition is used, the per-scene mean is subtracted. A fixed budget of primitives is selected per scene by opacity, then reordered along the Hilbert curve so that array index corresponds to a spatially stable location, which is what makes the slot-aligned reconstruction target learnable.

% ---------------------------------------------------------------------
\section{Gauge-Invariant Covariance Objective}
% ---------------------------------------------------------------------

The covariance objective compares the Gaussians through their covariance matrices rather than through the raw scale and quaternion, which removes the unlearnable orientation gauge described in \autoref{sec:method_overview}. For a predicted Gaussian with covariance $\Sigma_p$ and a target with covariance $\Sigma_t$, the shape discrepancy is the squared Bures distance
\begin{equation}
  d_{\mathrm{B}}^2(\Sigma_p, \Sigma_t) = \mathrm{tr}\!\left(\Sigma_p + \Sigma_t - 2\big(\Sigma_t^{1/2}\Sigma_p\,\Sigma_t^{1/2}\big)^{1/2}\right),
\end{equation}
which is the squared 2-Wasserstein distance between zero-mean Gaussians \citep{villani2009ot}. In practice we use the stable co-diagonalization surrogate $\lVert \Sigma_p^{1/2} - \Sigma_t^{1/2}\rVert_F^2$, which equals the Bures distance when the two covariances commute and is cheaper and better conditioned at coincident eigenvalues. Symmetric positive-definite square roots are computed by a trace-normalized Newton-Schulz iteration, which is differentiable and stable even at degenerate eigenvalues, with all covariance arithmetic carried out in single precision. The permutation-invariant set loss folds the same covariance square root into a per-primitive feature, so that a single Euclidean distance in feature space equals the weighted sum of the position, colour, opacity, and Bures-shape costs, and the within-block assignment is then solved by entropic optimal transport \citep{cuturi2013sinkhorn} before the matched pairs are scored.

% ---------------------------------------------------------------------
\section{Second-Stage Configuration}
% ---------------------------------------------------------------------

The second-stage models are flow-matching transformers in the style of a diffusion transformer \citep{peebles2023dit}, trained with the rectified-flow objective \citep{lipman2023flowmatching,liu2023rectified} on the frozen Stage~1 latent. The Stage~1 encoder and decoder are held fixed during all second-stage training: their parameters are frozen and set to evaluation mode, and the latent is encoded under no-gradient, so only the transformers are optimised. Each transformer uses a width of $384$, a depth of $12$, and $12$ attention heads, which is on the order of thirty to forty million parameters per model.

Token position is encoded with a one-dimensional rotary embedding \citep{su2021roformer} for the geometry and completion models, whose tokens lie along the Hilbert curve, and a learned absolute embedding for the layout model, whose tokens are abstract; this follows the small ablation of \autoref{sec:res_stage2}. Generation uses the hierarchical pair of \autoref{sec:method_stage2}: a layout model transports noise to the layout tokens, and a geometry model transports noise to the geometry tokens conditioned on the layout tokens through the same cross-attention interface used in the Stage~1 conditioning decoder. Completion fixes the encoded observed tokens and denoises the masked tokens; in the reported setup the observed and masked tokens are obtained by masking a contiguous block coverage of the latent of a fully encoded scene, so that each masked token corresponds to a spatial region. Because each model trains by encoding scenes on the fly rather than from a cache of precomputed latents, the second stage reuses the same training mixture as Stage~1 so that the frozen encoder is applied in-distribution. The three second-stage settings of \autoref{sec:res_selection}, namely generation and completion on the two disentangled checkpoints and completion without conditioning on the structured-flat checkpoint, share this configuration and differ only in which Stage~1 latent they consume and whether a layout model is present.



Reconstruction is reported as the parameter-space error, both as an aggregate and broken down by attribute, using the same reduction in training and validation so that the two are on the same scale. Because the raw quaternion error is blind to the orientation gauge, orientation and shape are additionally reported through a gauge-invariant covariance error, namely the mean per-primitive Bures and log-Euclidean distance between predicted and target covariances, together with the mean anisotropy of the target Gaussians, which indicates how much orientation is even constrained. Render fidelity is reported with peak signal-to-noise ratio, structural similarity \citep{wang2004ssim}, and the learned perceptual metric LPIPS \citep{zhang2018lpips} on images rendered from the predicted and ground-truth Gaussians. Generalization is monitored by the ratio of held-out to in-distribution error, where a ratio near one indicates no overfitting. Semantic structure is assessed qualitatively by principal-component visualizations of the per-primitive features and quantitatively by between-class versus within-class separation and a linear probe.


\end{document}

%%%%%%%%%%%%%%%%%%%%%%%%%%%%%%%%%%%%%%%%%%%%%%%%%%%%%%%%%%%%%%%%%%%%%%%%%%%%%%%%
%%%%%%%%%%%%%%%%%%%%%%%%%%%%%%%%%%%%%%%%%%%%%%%%%%%%%%%%%%%%%%%%%%%%%%%%%%%%%%%%